\documentclass[12pt,a4paper]{article}

% --- Pakiety ---
\usepackage[T1]{fontenc}
\usepackage[utf8]{inputenc}
\usepackage[polish]{babel}
\usepackage{geometry}
\usepackage{hyperref}
\usepackage{booktabs}
\usepackage{longtable}
\usepackage{array}
\usepackage{enumitem}
\usepackage{titlesec}
\usepackage{xcolor}
\usepackage{fancyhdr}
\usepackage{parskip}

% --- Marginesy ---
\geometry{
  a4paper,
  top=2.5cm,
  bottom=2.5cm,
  left=2.5cm,
  right=2.5cm
}

% --- Kolory ---
\definecolor{headerblue}{RGB}{31,73,125}
\definecolor{lightgray}{RGB}{245,245,245}
\definecolor{codegray}{RGB}{100,100,100}

% --- Nagłówki sekcji ---
\titleformat{\section}
  {\large\bfseries\color{headerblue}}
  {\thesection.}{0.6em}{}[\titlerule]

\titleformat{\subsection}
  {\normalsize\bfseries\color{headerblue}}
  {\thesubsection.}{0.5em}{}

\titleformat{\subsubsection}
  {\normalsize\bfseries}
  {\thesubsubsection.}{0.4em}{}

% --- Nagłówek i stopka ---
\pagestyle{fancy}
\fancyhf{}
\fancyhead[L]{\small\textit{ArcadePortal --- Dokumentacja projektu}}
\fancyhead[R]{\small\thepage}
\renewcommand{\headrulewidth}{0.4pt}

% --- Hiperłącza ---
\hypersetup{
  colorlinks=true,
  linkcolor=headerblue,
  urlcolor=headerblue,
  pdfauthor={Przemysław Pyrzyński, Michał Napora, Maksymilian Wojtkowski},
  pdftitle={ArcadePortal — Dokumentacja projektu}
}

% --- Czcionka kodowania (monospace) ---
\newcommand{\code}[1]{\texttt{\small#1}}

% ============================================================
\begin{document}

% --- Strona tytułowa ---
\begin{titlepage}
  \centering
  \vspace*{2cm}
  {\Huge\bfseries\color{headerblue} ArcadePortal\par}
  \vspace{0.5cm}
  {\Large Portal z grami przeglądarkowymi\par}
  \vspace{0.3cm}
  {\large Dokumentacja projektowa\par}
  \vspace{2cm}
  \rule{\textwidth}{0.4pt}
  \vspace{0.5cm}
  {\normalsize Projektowanie i programowanie systemów internetowych I\par}
  {\normalsize Semestr letni 2024/2025\par}
  \vspace{0.5cm}
  \rule{\textwidth}{0.4pt}
  \vspace{1.5cm}
  {\large Przemysław Pyrzyński \quad Michał Napora \quad Maksymilian Wojtkowski\par}
  \vfill
\end{titlepage}

% --- Spis treści ---
\tableofcontents
\newpage

% ============================================================
\section{Opis funkcjonalny systemu}

ArcadePortal to portal z grami przeglądarkowymi umożliwiający użytkownikom rejestrację kont,
rywalizację w grach zaimplementowanych w JavaScript oraz śledzenie wyników na globalnym rankingu.

\subsection{Moduł gier}

System obsługuje następujące gry przeglądarkowe:

\begin{itemize}[leftmargin=1.5cm]
  \item \textbf{Snake} --- klasyczna gra węża na siatce 24×24, rysowana na elemencie canvas
  \item \textbf{Tetris} --- zaawansowany Tetris z grawitacją, przyspieszeniem i podglądem następnego klocka
  \item \textbf{Flappy Bird} --- gra zręcznościowa z fizyką grawitacji i generowaniem rur
  \item \textbf{Pac-Man} --- klasyczny Pac-Man z duchami, power pelletami i systemem poziomów
  \item \textbf{Blackjack} --- karciana gra blackjack z pełną mechaniką asów i krupiera
\end{itemize}

\subsection{System kont użytkowników}

\begin{itemize}[leftmargin=1.5cm]
  \item Rejestracja z walidacją danych i emailem powitalnym
  \item Logowanie przez nazwę użytkownika lub adres email
  \item Profil użytkownika z historią rozgrywek i statystykami
  \item Zmiana nazwy użytkownika, avatara i hasła
  \item Reset hasła przez email z tokenem jednorazowym
  \item System ról: \textit{Admin} i \textit{User}
\end{itemize}

\subsection{Ranking i osiągnięcia}

\begin{itemize}[leftmargin=1.5cm]
  \item Globalny ranking oparty na sumie najlepszych wyników per gra
  \item Top 10 per gra z cachowaniem wyników (60 sekund)
  \item System odznak przyznawanych automatycznie po spełnieniu warunków
  \item Powiadomienia toast o nowych odznakach po zakończeniu gry
\end{itemize}

\subsection{Forum}

\begin{itemize}[leftmargin=1.5cm]
  \item Wątki ogólne i przypisane do konkretnej gry
  \item Tworzenie wątków, dodawanie odpowiedzi, edycja i usuwanie własnych postów
  \item Niezalogowani mogą czytać, zalogowani mogą pisać
  \item Moderacja przez administratorów
\end{itemize}

\subsection{Panel administratora}

\begin{itemize}[leftmargin=1.5cm]
  \item Zarządzanie użytkownikami: zmiana ról, blokowanie kont
  \item Zarządzanie grami: dodawanie, edycja, włączanie/wyłączanie
  \item Zarządzanie odznakami: tworzenie warunków per gra
  \item Usuwanie postów innych użytkowników
  \item Statystyki globalne: liczba użytkowników, sesji, wątków, postów
\end{itemize}

% ============================================================
\section{Opis technologiczny}

\subsection{Stos technologiczny}

\begin{longtable}{>{\bfseries}p{3cm} p{4cm} p{7cm}}
  \toprule
  Warstwa & Technologia & Opis \\
  \midrule
  \endhead
  Backend      & ASP.NET Core 8 MVC         & Framework MVC --- kontrolery, routing, autoryzacja \\[4pt]
  ORM          & Entity Framework Core 8     & Mapowanie obiektowo-relacyjne, migracje, LINQ \\[4pt]
  Baza danych  & SQLite                      & Lekka baza plikowa, idealna do projektu zaliczeniowego \\[4pt]
  Identity     & ASP.NET Core Identity       & System ról, hashowanie haseł, lockout, tokeny \\[4pt]
  Cache        & IMemoryCache                & Cachowanie rankingów w pamięci serwera \\[4pt]
  Frontend     & HTML5 + CSS3               & Razor Views, responsywny layout (RWD) \\[4pt]
  CSS Framework& Bootstrap 5                 & Grid system, komponenty UI \\[4pt]
  JavaScript   & Vanilla JS (ES6+)           & Logika gier, fetch API, animacje \\[4pt]
  Gry          & HTML5 Canvas API            & Renderowanie gier: Snake, Tetris, Flappy, Pac-Man \\[4pt]
  Mailing      & MailKit + SMTP              & Wysyłka maili: rejestracja, reset hasła, odznaki \\[4pt]
  Lokalizacja  & IStringLocalizer / IViewLocalizer & Obsługa języków PL i EN \\[4pt]
  Konteneryzacja & Docker + Docker Compose   & Pakowanie aplikacji z Mailpitem \\
  \bottomrule
\end{longtable}

\subsection{Architektura aplikacji}

Aplikacja zbudowana jest w architekturze MVC (Model-View-Controller) z dodatkową warstwą serwisów:

\begin{itemize}[leftmargin=1.5cm]
  \item \textbf{Models} --- encje EF Core mapowane na tabele bazy danych
  \item \textbf{DTOs} --- obiekty transferu danych między kontrolerami a widokami
  \item \textbf{Controllers} --- obsługa requestów HTTP, walidacja, logika routingu
  \item \textbf{Services} --- logika biznesowa (\code{AchievementService}, \code{MailKitEmailService})
  \item \textbf{Views} --- szablony Razor (\code{.cshtml}) z kodem HTML/CSS
  \item \textbf{Data} --- \code{AppDbContext}, \code{DbSeeder}, migracje EF Core
\end{itemize}

% ============================================================
\section{Schemat bazy danych (ERD)}

Baza danych oparta na SQLite zawiera 7 tabel własnych oraz tabele zarządzane przez
ASP.NET Core Identity (\code{AspNetUsers}, \code{AspNetRoles}, \code{AspNetUserRoles} i inne).

\subsection{Tabele własne}

\subsubsection{Users (rozszerzenie IdentityUser)}

\begin{longtable}{>{\ttfamily}p{3.5cm} p{3cm} p{7cm}}
  \toprule
  \normalfont\textbf{Kolumna} & \textbf{Typ} & \textbf{Opis} \\
  \midrule
  \endhead
  Id        & PK string   & Klucz główny (GUID generowany przez Identity) \\[3pt]
  UserName  & string      & Unikalna nazwa użytkownika \\[3pt]
  Email     & string      & Adres email użytkownika \\[3pt]
  PasswordHash & string   & Hash hasła (PBKDF2, zarządzany przez Identity) \\[3pt]
  AvatarUrl & string?     & Opcjonalny URL avatara użytkownika \\[3pt]
  CreatedAt & DateTime    & Data i czas rejestracji (UTC) \\
  \bottomrule
\end{longtable}

\subsubsection{Games}

\begin{longtable}{>{\ttfamily}p{3.5cm} p{3cm} p{7cm}}
  \toprule
  \normalfont\textbf{Kolumna} & \textbf{Typ} & \textbf{Opis} \\
  \midrule
  \endhead
  Id           & PK int    & Klucz główny \\[3pt]
  Slug         & string    & Unikalny identyfikator tekstowy (np. \texttt{"snake"}) \\[3pt]
  Title        & string    & Tytuł gry wyświetlany użytkownikom \\[3pt]
  Description  & string?   & Opcjonalny opis gry \\[3pt]
  ThumbnailUrl & string?   & URL miniaturki gry \\[3pt]
  IsActive     & bool      & Czy gra jest widoczna dla graczy \\[3pt]
  CreatedAt    & DateTime  & Data dodania gry (UTC) \\
  \bottomrule
\end{longtable}

\subsubsection{GameSessions}

\begin{longtable}{>{\ttfamily}p{3.5cm} p{3cm} p{7cm}}
  \toprule
  \normalfont\textbf{Kolumna} & \textbf{Typ} & \textbf{Opis} \\
  \midrule
  \endhead
  Id              & PK int    & Klucz główny \\[3pt]
  UserId          & FK string & Referencja do Users \\[3pt]
  GameId          & FK int    & Referencja do Games \\[3pt]
  Score           & int       & Wynik zdobyty w tej sesji \\[3pt]
  DurationSeconds & int       & Czas trwania rozgrywki w sekundach \\[3pt]
  PlayedAt        & DateTime  & Data i czas rozgrywki (UTC) \\
  \bottomrule
\end{longtable}

\subsubsection{Achievements}

\begin{longtable}{>{\ttfamily}p{3.5cm} p{3cm} p{7cm}}
  \toprule
  \normalfont\textbf{Kolumna} & \textbf{Typ} & \textbf{Opis} \\
  \midrule
  \endhead
  Id          & PK int    & Klucz główny \\[3pt]
  GameId      & FK int    & Referencja do Games (odznaka należy do gry) \\[3pt]
  Name        & string    & Nazwa odznaki wyświetlana użytkownikowi \\[3pt]
  Description & string    & Opis warunku zdobycia odznaki \\[3pt]
  IconUrl     & string?   & Emoji lub URL ikony odznaki \\[3pt]
  Condition   & string    & Warunek w formacie \texttt{"score>=100"} lub \texttt{"duration>=60"} \\
  \bottomrule
\end{longtable}

\subsubsection{UserAchievements}

\begin{longtable}{>{\ttfamily}p{3.5cm} p{3cm} p{7cm}}
  \toprule
  \normalfont\textbf{Kolumna} & \textbf{Typ} & \textbf{Opis} \\
  \midrule
  \endhead
  Id            & PK int    & Klucz główny \\[3pt]
  UserId        & FK string & Referencja do Users \\[3pt]
  AchievementId & FK int    & Referencja do Achievements \\[3pt]
  UnlockedAt    & DateTime  & Data i czas odblokowania odznaki (UTC) \\
  \bottomrule
\end{longtable}

\subsubsection{Threads}

\begin{longtable}{>{\ttfamily}p{3.5cm} p{3cm} p{7cm}}
  \toprule
  \normalfont\textbf{Kolumna} & \textbf{Typ} & \textbf{Opis} \\
  \midrule
  \endhead
  Id        & PK int    & Klucz główny \\[3pt]
  Title     & string    & Tytuł wątku (max 200 znaków) \\[3pt]
  UserId    & FK string & Autor wątku --- referencja do Users \\[3pt]
  GameId    & FK int?   & Opcjonalna referencja do Games (null = wątek ogólny) \\[3pt]
  IsPinned  & bool      & Czy wątek jest przypięty przez moderatora \\[3pt]
  CreatedAt & DateTime  & Data i czas utworzenia wątku (UTC) \\
  \bottomrule
\end{longtable}

\subsubsection{Posts}

\begin{longtable}{>{\ttfamily}p{3.5cm} p{3cm} p{7cm}}
  \toprule
  \normalfont\textbf{Kolumna} & \textbf{Typ} & \textbf{Opis} \\
  \midrule
  \endhead
  Id        & PK int    & Klucz główny \\[3pt]
  ThreadId  & FK int    & Referencja do Threads \\[3pt]
  UserId    & FK string & Autor posta --- referencja do Users \\[3pt]
  Body      & string    & Treść posta (max 5000 znaków) \\[3pt]
  CreatedAt & DateTime  & Data i czas utworzenia posta (UTC) \\[3pt]
  UpdatedAt & DateTime? & Data ostatniej edycji (null = nie edytowano) \\
  \bottomrule
\end{longtable}

\subsection{Relacje między tabelami}

\begin{longtable}{p{3.5cm} p{2.5cm} p{3.5cm} p{5cm}}
  \toprule
  \textbf{Tabela A} & \textbf{Relacja} & \textbf{Tabela B} & \textbf{Opis} \\
  \midrule
  \endhead
  Users        & 1 : N    & GameSessions    & Jeden użytkownik może mieć wiele sesji gry \\[3pt]
  Games        & 1 : N    & GameSessions    & Jedna gra może mieć wiele sesji \\[3pt]
  Users        & 1 : N    & UserAchievements & Jeden użytkownik może mieć wiele odznak \\[3pt]
  Achievements & 1 : N    & UserAchievements & Jedna odznaka może być odblokowana przez wielu \\[3pt]
  Games        & 1 : N    & Achievements    & Jedna gra definiuje wiele odznak \\[3pt]
  Users        & 1 : N    & Threads         & Jeden użytkownik może tworzyć wiele wątków \\[3pt]
  Games        & 0..1 : N & Threads         & Wątek opcjonalnie przypisany do gry \\[3pt]
  Threads      & 1 : N    & Posts           & Jeden wątek zawiera wiele postów \\[3pt]
  Users        & 1 : N    & Posts           & Jeden użytkownik może pisać wiele postów \\
  \bottomrule
\end{longtable}

% ============================================================
\section{Wdrożone zagadnienia kwalifikacyjne}

\begin{longtable}{>{\centering}p{0.7cm} p{3.5cm} p{9.5cm}}
  \toprule
  \textbf{\#} & \textbf{Zagadnienie} & \textbf{Implementacja} \\
  \midrule
  \endhead
  1  & Framework MVC           & ASP.NET Core 8 MVC --- kontrolery, modele, widoki Razor \\[3pt]
  2  & Framework CSS           & Bootstrap 5 --- grid, komponenty; własny system CSS z variables \\[3pt]
  3  & Baza danych             & SQLite przez Entity Framework Core 8 \\[3pt]
  4  & Cache                   & IMemoryCache --- cachowanie rankingów i leaderboardów \\[3pt]
  5  & Dependency manager      & NuGet --- zarządzanie pakietami .NET \\[3pt]
  6  & HTML                    & Razor Views (\code{.cshtml}) --- szkielet aplikacji \\[3pt]
  7  & CSS                     & Własne pliki CSS per moduł (auth, forum, game, admin...) \\[3pt]
  8  & JavaScript              & Logika gier w Vanilla JS, sterowanie, pętla gry \\[3pt]
  9  & Routing                 & Atrybutowy i konwencjonalny routing MVC, pretty URLs \\[3pt]
  10 & ORM                     & Entity Framework Core --- DbContext, migracje, LINQ \\[3pt]
  11 & Uwierzytelnianie        & ASP.NET Core Identity --- role, hasła, lockout, tokeny \\[3pt]
  12 & Lokalizacja             & IStringLocalizer / IViewLocalizer, pliki .resx, PL/EN \\[3pt]
  13 & Mailing                 & MailKit + SMTP --- rejestracja, reset hasła, odznaki \\[3pt]
  14 & Formularze              & Razor Forms z walidacją, AntiForgeryToken \\[3pt]
  15 & Asynchroniczne interakcje & Fetch API --- zapis wyników, asynchroniczne zapytania \\[3pt]
  16 & Konsumpcja API          & Własne REST API \code{/api/scores} konsumowane przez JS gier \\[3pt]
  17 & Publikacja API          & Endpoint \code{/api/scores} (POST) z autoryzacją i CSRF \\[3pt]
  18 & RWD                     & Bootstrap Grid + media queries, responsywny layout \\[3pt]
  19 & Logger                  & \code{ILogger<T>} we wszystkich kontrolerach i serwisach \\[3pt]
  20 & Deployment              & Docker + Docker Compose z Mailpit i volume na bazę \\
  \bottomrule
\end{longtable}

% ============================================================
\section{Instrukcja uruchomienia}

\subsection{Uruchomienie lokalne}

\textbf{Wymagania:} .NET 8 SDK, Mailpit

\begin{enumerate}[leftmargin=1.5cm]
  \item Sklonuj repozytorium: git clone \url{https://github.com/Przemek738/SystemyInternetowe_Projekt}
  \item Zainstaluj Mailpit i uruchom: \code{mailpit.exe}
  \item W katalogu projektu: \code{dotnet ef database update}
  \item Uruchom aplikację: \code{dotnet run}
  \item Otwórz przeglądarkę: \url{http://localhost:5000}
  \item Skrzynka mailowa: \url{http://localhost:8025}
\end{enumerate}

\subsection{Uruchomienie przez Docker}

\textbf{Wymagania:} Docker Desktop

\begin{enumerate}[leftmargin=1.5cm]
  \item Sklonuj repozytorium i przejdź do katalogu projektu
  \item Uruchom: \code{docker compose up --build}
  \item Aplikacja dostępna pod: \url{http://localhost:8080}
  \item Skrzynka mailowa: \url{http://localhost:8025}
\end{enumerate}

Przy pierwszym uruchomieniu seeder automatycznie tworzy role, konto admina, gry z odznakami
oraz 10 testowych użytkowników z losowymi wynikami.

% ============================================================
\section{Wnioski projektowe}

\subsection{Osiągnięte cele}

\begin{itemize}[leftmargin=1.5cm]
  \item Zrealizowano portal z 5 grami przeglądarkowymi w pełni zintegrowanymi z systemem zapisywania wyników
  \item Zaimplementowano kompletny system kont z rolami, uwierzytelnianiem i resetem hasła przez email
  \item Stworzono forum dyskusyjne z moderacją i wątkami przypisanymi do gier
  \item Skonteneryzowano aplikację przez Docker z pełną izolacją środowiska
\end{itemize}

\subsection{Napotkane trudności}

\begin{itemize}[leftmargin=1.5cm]
  \item Integracja własnego systemu menu gier z jednym widokiem \code{Play.cshtml} wymagała przemyślenia architektury JS
  \item Konfiguracja ASP.NET Core Identity z własnym modelem User wymagała znajomości cyklu życia tokenów i claims
\end{itemize}

\subsection{Możliwe rozszerzenia}

\begin{itemize}[leftmargin=1.5cm]
  \item Dodanie OAuth2 (logowanie przez Google/GitHub)
  \item System turniejów i rozgrywek wieloosobowych
  \item Migracja bazy danych z SQLite na PostgreSQL dla środowiska produkcyjnego
\end{itemize}

% ============================================================
\newpage
\section*{Opis zagadnień kwalifikacyjnych}
\addcontentsline{toc}{section}{Opis zagadnień kwalifikacyjnych}

W tej sekcji opisujemy szczegółowo każde z 20 zagadnień kwalifikacyjnych --- gdzie jest
zaimplementowane w projekcie oraz jak można sprawdzić że działa.

\begin{longtable}{>{\centering}p{0.7cm} p{2.8cm} p{4.5cm} p{5.5cm}}
  \toprule
  \textbf{\#} & \textbf{Zagadnienie} & \textbf{Gdzie w projekcie} & \textbf{Jak sprawdzić} \\
  \midrule
  \endhead
  1  & Framework MVC
     & \code{Controllers/*.cs}, \code{Program.cs}, \code{Views/*.cshtml}
     & Wejdź na \code{/Home}, \code{/Game}, \code{/Forum} --- każda strona to inny kontroler \\[4pt]
  2  & Framework CSS
     & \code{Views/Shared/\_Layout.cshtml} (Bootstrap CDN), \code{wwwroot/css/*.css}
     & Otwórz DevTools → Elements → sprawdź klasy Bootstrap i własne \\[4pt]
  3  & Baza danych
     & \code{Data/AppDbContext.cs}, \code{appsettings.json}, \code{arcade.db}
     & Otwórz plik \code{arcade.db} w DB Viewerze w Visual Studio \\[4pt]
  4  & Cache
     & \code{Controllers/GameController.cs} (\code{GetOrCreateAsync}), \code{LeaderboardController.cs}
     & Zagraj grę, sprawdź wynik --- przez 60s top 10 się nie zmienia \\[4pt]
  5  & Dependency manager
     & \code{ArcadeProject.csproj} (\code{<PackageReference>}), \code{Dockerfile}
     & Otwórz \code{.csproj} --- widać wszystkie paczki NuGet \\[4pt]
  6  & HTML
     & \code{Views/**/*.cshtml} --- każdy widok to szablon HTML
     & Ctrl+U na dowolnej stronie --- widać wygenerowany HTML \\[4pt]
  7  & CSS
     & \code{wwwroot/css/gamePlay.css}, \code{auth.css}, \code{forum.css}, \code{admin.css}, \code{leaderboard.css}, \code{profile.css}
     & DevTools → Network → CSS --- każda strona ładuje swój plik CSS \\[4pt]
  8  & JavaScript
     & \code{wwwroot/js/gamePlay.js}, \code{snake.js}, \code{tetris.js}, \code{flappy.js}, \code{pacman.js}, \code{blackjack.js}
     & DevTools → Sources → \code{wwwroot/js/} \\[4pt]
  9  & Routing
     & \code{Program.cs} (\code{MapControllerRoute}), \code{Controllers/GameController.cs} (\code{[HttpPost("/api/scores")]})
     & Przejdź na \code{/Game/Play/snake} --- \texttt{"snake"} to parametr z URL \\[4pt]
  10 & ORM
     & \code{Data/AppDbContext.cs}, \code{Models/*.cs}, \code{Migrations/}
     & Folder \code{Migrations/} --- każda migracja to zmiana w bazie \\[4pt]
  11 & Uwierzytelnianie
     & \code{Controllers/AccountController.cs}, \code{Program.cs} (\code{AddIdentity}), \code{Data/DbSeeder.cs}
     & Wejdź na \code{/Admin} bez logowania --- przekieruje na \code{/Account/Login} \\[4pt]
  12 & Lokalizacja
     & \code{Controllers/CultureController.cs}, \code{Resources/**/*.resx}
     & Kliknij EN w navbarze --- strona główna zmienia się na angielski \\[4pt]
  13 & Mailing
     & \code{Services/MailKitEmailService.cs}, \code{Services/IEmailService.cs}
     & Zarejestruj konto → sprawdź \url{http://localhost:8025} \\[4pt]
  14 & Formularze
     & \code{Views/Account/Register.cshtml}, \code{Views/Forum/Create.cshtml}, \code{DTOs/AccountDtos.cs}
     & Spróbuj zarejestrować się z pustymi polami --- pojawią się błędy \\[4pt]
  15 & Async interakcje
     & \code{wwwroot/js/gamePlay.js} (funkcja \code{saveScore()} z \code{fetch()})
     & DevTools → Network → zagraj grę → widać POST \code{/api/scores} \\[4pt]
  16 & Konsumpcja API
     & \code{wwwroot/js/gamePlay.js} (\code{fetch} POST \code{/api/scores})
     & Jak wyżej --- JS konsumuje endpoint wystawiony przez kontroler \\[4pt]
  17 & Publikacja API
     & \code{Controllers/GameController.cs} (metoda \code{SaveScore}, \code{[HttpPost("/api/scores")]})
     & DevTools → kliknij request \code{/api/scores} → Preview → JSON \\[4pt]
  18 & RWD
     & \code{wwwroot/css/gamePlay.css} (\code{@media max-width: 900px}), \code{\_Layout.cshtml} (\code{navbar-expand-md})
     & DevTools → Ctrl+Shift+I → symuluj iPhone → sprawdź układ \\[4pt]
  19 & Logger
     & \code{Controllers/HomeController.cs}, \code{GameController.cs}, \code{AccountController.cs}, \code{Services/MailKitEmailService.cs}
     & \code{docker compose logs app -f} → zaloguj się → widać logi \\[4pt]
  20 & Deployment
     & \code{Dockerfile}, \code{docker-compose.yml}, \code{.dockerignore}
     & Na innym komputerze: \code{docker compose up --build} → działa bez .NET \\
  \bottomrule
\end{longtable}

% ============================================================
\newpage
\section*{Opis szczegółowy każdego zagadnienia}
\addcontentsline{toc}{section}{Opis szczegółowy każdego zagadnienia}

\paragraph{1. Framework MVC}
Projekt używa ASP.NET Core 8 MVC. Każda strona to kontroler, który obsługuje żądania i zwraca widok.
Na przykład \code{HomeController} obsługuje stronę główną, \code{GameController} obsługuje listę gier
i uruchamianie gry, \code{ForumController} obsługuje forum itd. Routing jest skonfigurowany w \code{Program.cs}.

\paragraph{2. Framework CSS}
Używamy Bootstrap 5, który jest dołączony przez CDN w pliku \code{\_Layout.cshtml}. Bootstrap obsługuje
responsywny grid i nawigację. Oprócz Bootstrap mamy własne pliki CSS dla każdego modułu --- osobny plik
dla strony gry, forum, logowania, panelu admina itd. Dzięki temu style się nie mieszają i każda strona
ładuje tylko to, czego potrzebuje.

\paragraph{3. Baza danych}
Używamy SQLite, bo jest prosty w użyciu i nie wymaga osobnego serwera. Baza to jeden plik \code{arcade.db}.
W Dockerze ten plik jest trzymany na osobnym volume, żeby nie zginął przy rebuild.

\paragraph{4. Cache}
\code{IMemoryCache} jest wstrzyknięty do \code{GameController} i \code{LeaderboardController}. Ranking
top 10 per gra jest cachowany przez 60 sekund. Globalny ranking jest cachowany przez 2 minuty. Po zapisaniu
nowego wyniku cache jest usuwany przez \code{\_cache.Remove()}, żeby przy następnym wejściu dane były świeże.

\paragraph{5. Dependency manager}
NuGet jest standardowym menedżerem pakietów dla .NET. Wszystkie zależności (EF Core, Identity, MailKit,
Bootstrap itp.) są zapisane w pliku \code{ArcadeProject.csproj} w sekcji \code{PackageReference}. Komenda
\code{dotnet restore} pobiera je automatycznie. W Dockerfile restore jest jako osobna warstwa, co przyspiesza buildy.

\paragraph{6. HTML}
Używamy Razor Views --- pliki \code{.cshtml}, które są szablonami HTML z możliwością wstawiania kodu C\#
przez składnię \code{@}. Każdy kontroler ma swój folder w \code{Views/}. Plik \code{\_Layout.cshtml} jest
wspólnym szkieletem dla wszystkich stron --- zawiera navbar, footer i miejsca na sekcje Scripts i Styles.

\paragraph{7. CSS}
Oprócz Bootstrap mamy własne pliki CSS. Globalny styl jest w \code{site.css}. Każda sekcja ma swój plik ---
\code{auth.css} dla logowania i rejestracji, \code{forum.css} itd. Używamy CSS Custom Properties (zmienne)
dla kolorów motywu.

\paragraph{8. JavaScript}
Logika gier jest napisana w czystym JavaScript (ES6+) bez frameworków. Plik \code{gamePlay.js} to silnik
strony gry --- obsługuje menu, zapis wyników przez fetch i powiadomienia o odznakach. Każda gra ma swój plik
(\code{snake.js}, \code{tetris.js} itd.), który implementuje funkcję \code{window.gameStart()} wywoływaną
przez \code{gamePlay.js} po kliknięciu START. Komunikacja z serwerem odbywa się przez Fetch API.

\paragraph{9. Routing}
W \code{Program.cs} mamy standardowy routing MVC. Dzięki temu URL \code{/Game/Play/snake} przekazuje
\texttt{"snake"} jako parametr slug do metody \code{Play()} w \code{GameController}.

\paragraph{10. ORM}
Entity Framework Core 8 jest naszym ORM. \code{AppDbContext} dziedziczy po \code{IdentityDbContext}
i definiuje DbSety dla wszystkich encji. Relacje (np. kaskadowe usuwanie) są skonfigurowane w metodzie
\code{OnModelCreating}. Zapytania piszemy w LINQ. Migracje są w folderze \code{Migrations/}.

\paragraph{11. Uwierzytelnianie}
ASP.NET Core Identity zarządza użytkownikami i rolami. Mamy dwie role: \textit{Admin} i \textit{User}
tworzone przez seeder przy starcie. Hasła są hashowane algorytmem PBKDF2. Po 5 nieudanych próbach konto
jest blokowane na 5 minut. Atrybut \code{[Authorize(Roles="Admin")]} na \code{AdminController} sprawia,
że tylko admini mają dostęp do panelu. Logowanie działa przez username lub email.

\paragraph{12. Lokalizacja}
Obsługujemy dwa języki: polski i angielski. W navbarze jest przycisk PL/EN, który wysyła POST do
\code{CultureController} i zapisuje wybrany język w ciasteczku na rok. Tłumaczenia są w plikach \code{.resx}
w folderze \code{Resources/} --- osobny plik per kontroler i widok. W widokach używamy
\code{@inject IViewLocalizer Loc} i \code{@Loc["klucz"]} zamiast hardcodowanych tekstów.

\paragraph{13. Mailing}
Do wysyłania emaili używamy biblioteki MailKit. Stworzyliśmy interfejs \code{IEmailService} z metodami
\code{SendWelcomeAsync}, \code{SendPasswordResetAsync} i \code{SendAchievementAsync}. Implementacja
\code{MailKitEmailService} łączy się przez SMTP do Mailpit (lokalnie i w Dockerze) lub zewnętrznego serwera.
Jeśli wysyłka się nie uda, błąd jest logowany, ale nie przerywa działania aplikacji (try/catch).

\paragraph{14. Formularze}
Formularze w aplikacji używają Razor Tag Helpers --- \code{asp-for}, \code{asp-action},
\code{asp-validation-for}. Walidacja jest definiowana przez atrybuty w klasach DTO, np. \code{[Required]},
\code{[MinLength(6)]}, \code{[EmailAddress]}. Każdy formularz POST jest zabezpieczony przez
\code{[ValidateAntiForgeryToken]}, który chroni przed atakami CSRF. Błędy walidacji wracają do widoku
przez \code{ModelState}.

\paragraph{15. Asynchroniczne interakcje}
Po zakończeniu gry JavaScript wywołuje \code{fetch()} POST \code{/api/scores} bez przeładowania strony.
Wysyła JSON z wynikiem i czasem gry. Po odebraniu odpowiedzi aktualizuje pasek statusu (wynik, rekord osobisty)
i pokazuje powiadomienia o odznakach jako animowane elementy HTML dodawane do body. Wszystko dzieje się bez reload.

\paragraph{16. Konsumpcja API}
Plik \code{gamePlay.js} konsumuje endpoint \code{/api/scores} przez natywne \code{fetch()} API przeglądarki.
Wysyła JSON body z \code{gameSlug}, \code{score} i \code{durationSeconds}. Dołącza też token CSRF z ukrytego
pola formularza w nagłówku \code{RequestVerificationToken}. Obsługuje błędy --- jeśli serwer zwróci 401
(niezalogowany), wyświetla komunikat bez crashowania gry.

\paragraph{17. Publikacja API}
Metoda \code{SaveScore()} w \code{GameController} jest wystawiona jako REST endpoint. Ma atrybut
\code{[HttpPost("/api/scores")]} i \code{[ValidateAntiForgeryToken]}. Sprawdza czy użytkownik jest zalogowany
(401 jeśli nie) i czy gra istnieje (404 jeśli nie). Zapisuje sesję do bazy, wywołuje \code{AchievementService}
i zwraca JSON: \code{\{ message, score, newAchievements: [...] \}}.

\paragraph{18. RWD}
Strona jest responsywna dzięki Bootstrap Grid i własnym media queries. Strona gry używa CSS Grid, który
automatycznie przełącza się z 2 kolumn (gra + sidebar) na 1 kolumnę poniżej 900px. Nawigacja korzysta
z \code{navbar-expand-md} Bootstrap --- menu zwija się do hamburgera na telefonach. Siatka kart gier używa
\code{auto-fill} z \code{minmax(260px, 1fr)}, co automatycznie dobiera liczbę kolumn.

\paragraph{19. Logger}
Każdy kontroler i serwis ma wstrzyknięty \code{ILogger<T>} przez konstruktor. Logujemy informacje przy
kluczowych akcjach: załadowanie strony głównej, logowanie, rejestracja, wejście w grę, zapis wyniku,
przyznanie odznaki. Błędy wysyłki emaila logujemy przez \code{LogError}. W Dockerze logi widać przez
komendę \code{docker compose logs app -f}.

\paragraph{20. Deployment}
Aplikacja jest skonteneryzowana w Dockerze. \code{Dockerfile} używa multi-stage build --- najpierw obraz
z SDK do kompilacji, potem lżejszy obraz \code{aspnet} do uruchomienia. \code{docker-compose.yml} uruchamia
dwa kontenery: aplikację na porcie 8080 i Mailpit na portach 8025/1025. Baza SQLite jest na volume
\code{db-data}, który przeżywa rebuildy. Na innym komputerze wystarczy Docker Desktop i komenda
\code{docker compose up --build}.

\end{document}
